\section{The pole in the downward theta integral}

% BEGIN Y86N1
\paragraph{Y86N1. The partial Euler product at the original long string.}
Let $S$ be a finite set of finite places of $F$, and omit all places
of $E$ above it in every product below. Suppose the unramified
multiset of a rank-$(r+t)$ system, after division by a unitary
character $z_w$, has a rank-$r$ remainder $C_w$ and the $t$ entries
$q_w^{(t-1)/2-j}$ for $0\le j<t$, where $q_w=\mathrm Nw$.
Assume $|b|\le q_w^{(r-1)/2}$ for $b\in C_w$ and $r+5\le t$.
This includes the source's $\mathrm{SimplePole}$ hypothesis:
there $C_w$ is the multiset $c_w$ divided by its other norm-one
scalar $m_w$. Set
\[
 L_C^S(s)=\prod_{w\nmid S}\prod_{b\in C_w}(1-bq_w^{-s})^{-1},
 \qquad s_0=(t+1)/2 .
\]
On each compact subset of $\Re s>(r+1)/2$ the logarithmic series
$\sum_{w\nmid S}\sum_{b\in C_w}\sum_{h\ge1}b^h q_w^{-hs}/h$ converges
normally. Indeed $|bq_w^{-s}|\le q_w^{-1-\delta}$ for some
$\delta>0$ there, and $\sum_w q_w^{-1-\delta}<\infty$ by the
absolutely convergent Dedekind zeta product. Its exponential is
$L_C^S(s)$, holomorphic and nowhere zero on that half-plane. For
$\Re s>s_0$ all the string factors converge absolutely as well, so
the twisted standard Euler product is
\[
 L_C^S(s)\prod_{j=0}^{t-1}\zeta_E^S(s+(t-1)/2-j).
\]
The ordering reverses the root exponents, which preserves their
multiset. The partial Dedekind zeta function is its classical
meromorphic continuation times the omitted factors
$\prod_{w\mid S}(1-q_w^{-s})$, with nonzero residue at $1$.
This yields continuation of the displayed product to
$\Re s>(r+1)/2$. At $s_0$ the zeta arguments are $t,t-1,\ldots,1$.
Exactly the last factor has a pole, since $\zeta_E^S(h)>0$ for real
$h>1$. Its product residue is therefore
\[
 L_C^S(s_0)\,\operatorname*{Res}_{s=1}\zeta_E^S(s)
        \prod_{h=2}^t\zeta_E^S(h)\ne0 .
\]
Every continuation used in the source agrees with this one by
meromorphic uniqueness from the right half-plane. This proves the
full partial-Euler-product assertion $\mathrm{SimplePole}$.

For a native good point, take a good witness list and any upper
label $P$ above the point. R86G4 gives one lower cusp; H86P3
identifies its augmented base-change list with a list for $P$
up to permutation. At every place where the list is unramified
its distinguished length-$t$ string gives the entries above after
division by the Satake value of $\mu_i\eta^{-1}$. The
remaining cusp is twisted by unitary characters only. S86A2 gives
$|b|<q_w^{1/2}\le q_w^{(r-1)/2}$, since the member has $r\ge2$.
Choose $S$ to contain the places under all ramified factors.
Both the rank guard and the pole just proved apply to every such
$S$ on this member. With the $s+1/2$ argument of the doubling
integral, the pole in its variable is $s=t/2$.
% END Y86N1

% BEGIN Y86N2
\paragraph{Y86N2. A conditional descent from the residue to the prescribed test.}
Fix a nonzero vector $y$ at an original good upper point $x$.
Use the upper realization of L86M1, and write $I_y(f,\ell,g)$ for
the downward theta integral over the member's definite source
quotient, with Schwartz function $f$, coefficient functional $\ell$
and target argument $g$. The following sufficient hypotheses
specify which assertions must hold for this same $y$.

First, suppose there are a scalar component
$u(h)=\ell(\operatorname{ext}_V(y)(h))$ and a polynomial Fock test
$\Phi$ spherical outside a finite set $S$ and principal-level
invariant at the excluded places, with the source's nonzero
archimedean integral of the Weil and automorphic matrix coefficients
for $(\Phi,u)$. Suppose $u$ is a cuspidal principal-level form with
the normalized unramified parameters of Y86N1 outside $S$.
Use the member's Weil representation, characters and Haar measures.
Suppose the doubling integral $Z_y(s)$ formed from this pair after
adjoining $t$ hyperbolic planes admits meromorphic continuation from
the region of absolute convergence to a connected region containing
$s=t/2$. Require the unramified unfolding and the excluded-place
calculation there to give, initially on a right half-plane,
\[
 Z_y(s)=C_y\,L^S(s+1/2)/d^S(s),\qquad C_y\ne0.
\]
Here the numerator is precisely the twisted product of Y86N1.
The normalizer is the product over $0\le h<a$ of
$L_F^S(2s+a-h,\eta_{E/F}^{b+h})$, with $a=r+t$, $b=r+2t$ and
$\eta_{E/F}$ the quadratic character. When $\Re s>0$ each argument
has real part greater than $1$. Absolute convergence of the logarithms
makes every factor holomorphic and nonzero near $t/2$.
Equality by continuation is taken on the component reached from the
right within the common domain and containing $t/2$.
Y86N1 then gives $\operatorname{Res}_{s=t/2} Z_y(s)\ne0$.
Assume the regularized Siegel--Weil comparison for these same data
makes this residue zero whenever all $I_y(f,\ell',g)$ vanish.
Some test, component and target argument therefore give a nonzero
integral for $y$.

For descent to the prescribed test, assume the following local
generation properties in the member's actual realization.
On the good component the inclusion of the native downward
operator $\mathrm{thG}$ equals the full fixed-test operator
$\mathrm{thF}$. Vanishing of $\mathrm{thF}(y)$ annihilates the
integrals against the designated polynomial Fock tests. Their
archimedean generation and approximation, in a topology for which
each pairing $I_y(-,\ell',g)$ is continuous, give this vanishing
on every Schwartz function whose finite Heisenberg annihilator
contains that of the member's vacuum $v_P$. Here the annihilators
consist of finitely supported group-algebra sums of Heisenberg
translations with zero archimedean component.

Require the whole Schwartz space to be spanned, allowing closure
where these pairings are continuous, by target translates of that
class, source-level character differences, and the weighted
mirabolic averages from the rank-deficient finite orbits.
The level action on $u$ for each coefficient and right translation
of Haar measure kill the differences. Require that for every
designated prime and every nonzero defining vector, the corresponding
weighted subgroup averages annihilate every scalar component of $y$.
These are the subgroups fixing the indicated hyperplane pointwise
and the subgroups acting trivially modulo the indicated line,
with both determinant-character exponents $1$ and $-1$.
Their vanishing, transported through the integral, kills the
averages. Target translations only change $g$.
Consequently $\mathrm{thG}(y)=0$ would give $I_y(f,\ell',g)=0$
for every test, every coefficient and every $g$, contradicting the
nonzero residue. If the hypotheses hold for every nonzero vector in
the indicated good eigenspace, the native downward operator is
injective there. The occurrence of a varying nonzero test establishes
only the first implication. The factorization, residue comparison
and fixed-test hypotheses above remain inputs to this criterion.
% END Y86N2
